\documentclass[11pt,a4paper,openany]{book}
\input{userguide-preamble}
\hypersetup{pdftitle={HugMesh Games --- Full-stack developer guide}}
\begin{document}

\begin{titlepage}\thispagestyle{empty}
\noindent\begin{tikzpicture}[remember picture,overlay]
  \fill[snapBg] (current page.south west) rectangle (current page.north east);
  \fill[snapAccent,opacity=0.18]    ([yshift=-0.05\paperheight]current page.north west) circle (8cm);
  \fill[snapAccentTwo,opacity=0.10] ([xshift=4cm,yshift=-2cm]current page.north west) circle (5cm);
\end{tikzpicture}
\vspace*{4cm}\noindent
{\sffamily\bfseries
  {\color{snapAccentTwo}\fontsize{18}{22}\selectfont HUGMESH GAMES}\\[3mm]
  {\color{white}\fontsize{42}{50}\selectfont Full-stack\\developer guide}\par}
\vspace{1.5cm}\noindent\color{white!75!snapAccentTwo}\rule{0.4\linewidth}{1.2pt}
\vspace{0.8cm}\noindent
{\sffamily\color{white!85}\large
DB to browser, with the rules that hold the federated stack together.}
\vfill\noindent
{\sffamily\color{white!70}\small
Companion to \texttt{fullstack-developer.pptx}. Generated: \snapbuilddate{} \snapbuildtime}
\end{titlepage}

\frontmatter\pagestyle{empty}\tableofcontents\clearpage
\mainmatter\pagestyle{fancy}

\chapter{Mental model}
One repository, three languages, four runtimes. The BFF is Java 21 +
Spring Boot 4.0 + WebFlux + R2DBC + JDBC. The frontend is TypeScript +
Angular standalone components + Module Federation. The games are
GDScript + Godot 4.6, with a Rust GDExtension shared across them. The
GraalVM native AOT build of the BFF is the same Java code with a
sub-second cold start.

\chapter{How requests are shaped}
\snapdiagram{topology.tex}

The browser hits the host nginx; the host routes to the matching
per-remote nginx for \texttt{/\_remotes/<slug>/...} and to the BFF for
\texttt{/api/...}. Inside the BFF every controller is a
\texttt{@RestController}; routing is annotation-based, and WebFlux
decides reactor vs blocking by inspecting return types. Liquibase runs
once at boot on a JDBC pool of size 3 (separate from the R2DBC runtime
pool).

\chapter{Backend tour, in five files}
\begin{itemize}
  \item \snapfile{GameStreamController} --- the SSE merge that drives
        the live HUD. The \texttt{autoCancel=false} rule lives here.
  \item \snapfile{GameUploadService} --- the install transaction.
        Understanding this gives the whole hot-upload story.
  \item \snapfile{FederationController} --- the manifest endpoint that
        drives the host's import map. Trivial code, foundational
        consequence.
  \item \snapfile{JdbcDataSourceConfig} --- the one-line
        \texttt{@Primary HikariDataSource} that explains R2DBC + JDBC
        coexistence.
  \item \snapfile{R2dbcConvertersConfig} --- registers every entity for
        native-image reflection. Adding a new entity touches one place.
\end{itemize}

\chapter{Reactive vs blocking}
\snapdiagram{game-stream-controller.tex}

WebFlux requires non-blocking IO on the dispatcher thread. R2DBC for
the runtime path; JDBC + Liquibase only at boot. If a chain needs to
call something blocking, wrap it:

\begin{lstlisting}[language=Java]
Mono.fromCallable(() -> blockingCall())
    .subscribeOn(Schedulers.boundedElastic())
\end{lstlisting}

\begin{snapcaution}[title={Sinks.many() must set autoCancel=false}]
The default kills the sink on the first reconnect. Long-lived merges
go through \texttt{Flux.merge} --- see
\snapfile{GameStreamController} for the canonical 5-channel pattern.
\end{snapcaution}

\chapter{Adding a federated route}
\snapdiagram{federation-bootstrap.tex}

\begin{enumerate}
  \item Create \snapfile{frontend/projects/<remote>/} with its own
        \texttt{federation.config.js}.
  \item Expose the entry component in the \texttt{exposes} map.
  \item Insert a \texttt{federation\_modules} row (Liquibase
        changeset). \texttt{kind=remote-shell} for full route
        components.
  \item Add a \texttt{loadRemote(...)} entry to
        \snapfile{frontend/src/app/app.routes.ts}.
  \item Add a \texttt{/\_remotes/<remote>/} proxy block to
        \snapfile{frontend/nginx.conf}.
  \item Add the per-remote service to \snapfile{docker-compose.yml}.
  \item Redeclare DTOs inside the remote --- types erase across the
        federation boundary.
\end{enumerate}

\chapter{Common code-review patterns}
\begin{itemize}
  \item New endpoint? Check that
        \snapfile{R2dbcConvertersConfig} is updated and
        \texttt{make gen-api} ran.
  \item New SSE channel? Check
        \texttt{Sinks.many(...).onBackpressureBuffer(N, false)} ---
        the \texttt{false} is autoCancel.
  \item New federated remote? Check the host
        \texttt{nginx.conf}'s \texttt{/\_remotes/} block AND the
        \texttt{federation\_modules} row.
  \item New Liquibase changeset? Next integer prefix and included in
        \snapfile{db.changelog-master.xml}.
  \item New game? Four-place slug contract (see Godot guide).
\end{itemize}

\chapter{Per-session SSE sinks}
The single global sink that drove the original Three Wizards demo
became a leak as soon as two players ended up in different sessions:
every event landed on every subscriber. The current shape is a
\texttt{ConcurrentHashMap<String, SessionSinks>} keyed on
\texttt{sessionId}, lazy-allocated on first emit. The signature that
matters:

\begin{lstlisting}[language=Java]
applyControl(ControlEvent event,
             String characterId,
             String sessionId)
\end{lstlisting}

The legacy 2-arg form still exists for the single-player path; the
3-arg form is what every multiplayer-aware controller must call.
\texttt{GameStreamController} accepts an optional \texttt{?sessionId=}
query param and filters its subscription accordingly.

The isolation contract --- ``emit on session A's sinks; session B's
per-session subscriber sees nothing'' --- is locked in by
\snapfile{SessionIsolationTest} on the BFF side and
\reposrc{e2e/features/multiplayer/rpg-two-players.feature} on the
browser side. \snapfile{tests/cross-session-isolation.spec.ts} is the
explicit negative-direction spec.

Live admin observation goes through
\texttt{/api/admin/sessions/\{id\}/stream}, which merges movement +
world + character + view per-session sinks into one Flux. The
per-player \emph{view} channel is what makes it useful: HostBridge's
\texttt{view\_egress} posts $\sim$10 Hz camera samples to
\texttt{/api/view}, fanned out on the same per-session sink.

\chapter{Mod / publisher scope schema}
Roles live in Casdoor; per-game scope lives in Postgres. Two layers,
both required, neither sufficient alone:

\begin{itemize}
  \item Casdoor membership in \texttt{ROLE\_MOD} /
        \texttt{ROLE\_PUBLISHER} --- claims land in the JWT.
  \item One row in \texttt{mod\_game\_scopes} /
        \texttt{publisher\_game\_scopes} per (casdoorId, gameSlug)
        pair --- the BFF rejects any mod/publisher action whose
        target game isn't in the caller's active scope.
\end{itemize}

Adding a new scope check follows three steps:

\begin{enumerate}
  \item In the controller, gate the class with
        \texttt{@PreAuthorize("hasAnyRole('PUBLISHER','ADMIN')")} (or
        the mod equivalent).
  \item In the service, look up the row before mutating:
\begin{lstlisting}[language=Java,xleftmargin=4mm]
publisherGameScopeRepository
    .existsByCasdoorIdAndGameSlug(casdoorId, slug);
\end{lstlisting}
        Reject with 403 when admin-bypass doesn't apply.
  \item Add a Gherkin rejection scenario keyed to the new endpoint
        --- the format the existing
        \reposrc{e2e/features/security/mod-grant-gates.feature}
        uses (unauthenticated 401, plain player 403,
        wrong-scope publisher 403).
\end{enumerate}

Mutation endpoints follow a strict admin / publisher split:

\begin{lstlisting}[language=bash,xleftmargin=4mm]
# Admin-only: grant or revoke on any game.
/api/admin/role-grants/grant-mod
/api/admin/role-grants/revoke-mod
/api/admin/role-grants/grant-publisher
/api/admin/role-grants/revoke-publisher

# Publisher: grant mod, but only on a game in their scope.
/api/publisher/grant-mod
\end{lstlisting}

These are the only paths that add or remove scope rows. The publisher
endpoint walks the same write path as the admin one but enforces
caller-owns-target-game before mutating.

\chapter{Rate-limit layers}
Two layers of throttling, deliberately split:

\paragraph{Layer 1 --- nginx zones.}
\snapfile{frontend/nginx.conf} and \snapfile{frontend/nginx-tls.conf}
declare \texttt{limit\_req\_zone} entries sized for $\sim$10 000
concurrent users. Per-zone caps: \texttt{browse} 100 r/s,
\texttt{api} 50 r/s, \texttt{sse} 30 r/s (nodelay), \texttt{auth}
+ \texttt{admin} 5 r/s (nodelay), \texttt{maker\_api} 10 r/s. Plus
\texttt{limit\_conn conn\_per\_ip 40} for parallel-socket caps.
Tripped requests return 429, not 503, so well-behaved clients back
off. Cheap edge rejection before traffic ever touches the JVM.

\paragraph{Layer 2 --- in-pod + cross-pod.}
\snapfile{RateLimitFilter} runs the local in-process bucket as the
fast pre-check, then augments cluster-wide via
\snapfile{DistributedRateLimitService} --- an append-row log on the
Cassandra \texttt{rate\_limit\_events} table with a 90 s TTL, plus a
window-bounded count against
\texttt{event\_id $>=$ minTimeuuid(now $-$ WINDOW)}. Disable with
\texttt{app.rate-limit.distributed.enabled=false} (default
\texttt{true}); the filter falls back to local-only when Cassandra
is unreachable.

\begin{snapcaution}[title={Don't replace one layer with the other}]
The nginx zones reject at the edge before the JVM spins up a request
handler --- they exist so a flood can't exhaust the BFF's bounded-
elastic scheduler. The Java-side filter exists so a user can be
shaped per-user, not just per-IP, and so the cluster sees the same
budget regardless of which pod a request lands on. Removing either
opens a hole.
\end{snapcaution}

\chapter{Security primitives}
Three small APIs the codebase has converged on:

\begin{itemize}
  \item \snapfile{Crypto.mintUid()} --- the single mint for every
        UI-facing UID and URL token. Output: 43-character URL-safe
        base64 of \texttt{SHA3-256(64 random bytes from BouncyCastle
        SP800-90A HMAC-SHA256 DRBG)}. No more
        \texttt{UUID.randomUUID()} or raw \texttt{SecureRandom}.
        For raw bytes: \texttt{Crypto.randomBytes(n)}; for a hash:
        \texttt{Crypto.sha3\_256(...)}. The DRBG state is a
        process-singleton.
  \item \snapfile{Roles.has(jwt, "ADMIN")} /
        \texttt{Roles.fromAuthentication(auth)} /
        \texttt{Roles.isAdmin(jwt)} --- centralised role extraction.
        Don't copy the claim walk; use the helper. Controller-level
        access checks should use
        \texttt{@PreAuthorize("hasRole('ADMIN')")}; method-internal
        checks fall through to \texttt{Roles}.
  \item Liquibase 036 widens session-id columns from
        \texttt{varchar(36)} to \texttt{varchar(64)} to fit the new
        SHA3 shape (43 chars). Pre-existing UUID-shape rows continue
        to fit.
\end{itemize}

Session-scoped control authorization: two players sharing the same
\texttt{character\_id} in different sessions used to 403 each other
because \texttt{findActiveOwnerOfCharacter(character\_id)} raced on
Postgres' first-row pick. The fix is
\texttt{findActiveOwnerInSession(sessionId, characterId)} +
the 3-arg \texttt{authorizeControl(casdoorId, characterId,
sessionIdOpt)} overload; \texttt{ControlController} threads the
\texttt{X-Session-Id} header through. Pinned in
\snapfile{GameSessionServiceTest}.

\chapter{What NOT to do}
\begin{itemize}
  \item Don't call JDBC from a WebFlux controller.
  \item Don't import \snapfile{@angular-architects/native-federation}
        statically from a route handler.
  \item Don't share TypeScript DTO types across federation boundaries.
  \item Don't bypass Liquibase. Manual SQL on a live DB diverges.
  \item Don't add controller logic without tests AND an updated
        Javadoc.
\end{itemize}

\end{document}
